\section{Introduction}
\label{sec:introduction}

Wi-Fi and 5G New Radio in unlicensed spectrum (NR-U) contend without a common
scheduler. Both defer to detected energy and count down a backoff, yet their
surrounding procedures differ. Wi-Fi follows DCF/EDCA after an inter-frame
deferral. NR-U category-4 LBT also observes channel-occupancy limits and aligns
access to an NR transmission boundary
\cite{ieee80211_2020,3gpp37213,naik2020next}. Simultaneous counter expiry can
therefore waste airtime across both technologies.

DB-LBT replaces part of repeated random recovery with a counter derived from
locally observed countdown interruptions \cite{kosek2022dblbt}. The mechanism
can approach a collision-reduced order without node identities or
cross-technology messages. Its later TMC form remains effective under
non-full-buffer traffic, random busy periods, sensing perturbations, and
asymmetric energy detection \cite{tinnirello2026deterministic}. Those cases
also expose a control problem. A recovery tuple that admits intermittent
queues quickly can be too aggressive after more contenders become active,
whereas a conservative tuple can add avoidable delay in a light cell.

The instantaneous radio channel is not predictable from a short MAC history.
It includes nonlinear CCA threshold crossings, fading, interference, and
decoding outcomes. A transmitter nevertheless observes a slower protocol
trace: countdown freezes, queue growth, access delay, retransmissions, and
ACK/HARQ outcomes. When load or occupancy persists for several update windows,
that trace can identify which DB-LBT profile works better without reconstructing
the channel itself. This motivates our main formulation:
\emph{DB-LBT parameter adaptation is a protocol-constrained contextual bandit
under locally observable MAC information.}

The formulation targets mutually detectable indoor cells with piecewise-stable
offered load or active population. The context must be stable long enough to
learn, yet different enough across phases to make adaptation worthwhile. It is
not intended for per-packet channel prediction.

The contributions are threefold. First, we map DB-LBT recovery control to a
contextual bandit whose actions cannot alter CCA thresholds, Wi-Fi deferral,
NR-U synchronization or occupancy limits, PHY rate, or retransmission
semantics. The 24-profile design space is screened offline and reduced to two
conflicting profiles for online selection. Second, we connect the 11 local
features to protocol events and state explicit sensing, reception,
separability, and dwell-time conditions under which they are informative.
Third, a paired 32-seed experiment measures both performance and adaptation
time across repeated $4+4$ and $6+6$ phases. It finds a positive
phase-conditional effect, identifies a small high-load penalty, and preserves
the opposite whole-run diagnostic rather than hiding aggregation sensitivity.
